\BeleskeID{09_2-s034}
% element: 09_2-s034:d01
% element: 09_2-s034:e04
Zamenom veze između ulaznog i izlaznog napona u polaznu ravnotežu struja, za jednake parametre i suprotne pragove dobijamo
\[\left(V_O+\frac{V_{DD}}2-V_{DD}+V_{Tn}\right)^2=V_O\left(2\left(V_O+\frac{V_{DD}}2-V_{Tn}\right)-V_O\right)\]
\[\left(V_O-\frac{V_{DD}}2+V_{Tn}\right)^2=V_O(V_O+V_{DD}-2V_{Tn})\]
Razvijanjem kvadrata članovi $V_O^2$ se poništavaju. Preostala jednačina linearna je po $V_O$, pa daje izlazni napon u tački gornjeg praga. Ulazni napon se potom dobija iz $V_I=V_O+V_{DD}/2$.
